
\section{Il generatore di carico e le opzioni di strumentazione}

Il workload \`e generato da rt-app, un programma che esegue task periodici descritti in un file di
configurazione. La versione utilizzata \`e stata fornita come base di partenza, gi\`a strumentata
con OpenTelemetry, e ne abbiamo studiato il funzionamento per poterla configurare correttamente.

\subsection{Come si descrive un workload}

La configurazione \`e un file JSON con due sezioni. La sezione globale stabilisce la durata
dell'esecuzione, la costante di calibrazione e la destinazione dei log; la sezione dei task
descrive, per ciascun task, la politica di scheduling, la priorit\`a, le CPU su cui pu\`o girare e
la sequenza di eventi che ne costituisce la fase periodica.

\begin{lstlisting}[caption={Struttura della configurazione: un task critico e un task di disturbo.}]
{
  "global": {
      "duration": 20,          // durata dell'esecuzione, in secondi
      "calibration": 29,       // costo per iterazione, fissato (vedi 1.5)
      "logdir": "...",         // cartella dei log, una per esecuzione
      "log_basename": "rtapp"
  },
  "tasks": {
      "HI_task": {                    // task critico
          "policy": "SCHED_FIFO",     // scheduling real-time
          "priority": 90,
          "cpus": [2],                // core fisico 1
          "loop": -1,                 // ripete fino a fine durata
          "run": 2000,                // 2000 us di lavoro effettivo
          "timer": { "ref": "unique", "period": 10000, "mode": "absolute" }
      },
      "LO_noise": {                   // task di disturbo
          "policy": "SCHED_OTHER",
          "instance": 4,              // repliche indipendenti
          "cpus": [6],                // core fisico 3
          "loop": -1,
          "run": 1000,
          "sleep": 1000
      }
  }
}
\end{lstlisting}

Gli eventi che compongono la fase sono di tre tipi: \texttt{run} esegue un ciclo di calcolo per la
durata indicata; \texttt{sleep} sospende il task per un tempo relativo alla fine del lavoro;
\texttt{timer} sospende il task fino a un istante calcolato su una griglia temporale.

\textbf{Il task critico usa un timer su griglia assoluta}, e questa \`e una scelta deliberata. Con
un'attesa relativa, l'istante di attivazione successivo deriva dalla somma fra lavoro e attesa, e
ne eredita quindi tutta la variabilit\`a; inoltre, in questa modalit\`a il programma non calcola il
margine residuo rispetto alla scadenza, che \`e la grandezza su cui si fonda tutta l'analisi
successiva. Il confronto diretto fra le due modalit\`a, a parit\`a di tutto il resto, \`e riportato
in Tabella~\ref{tab:timer}.

\begin{table}[h]\centering\small
\begin{tabular}{lcccc}
\toprule
 & iterazioni & periodo mediano & jitter & margine calcolato \\
\midrule
attesa relativa & 1981/2000 & 10062\us & 37\us & \textbf{mai} \\
timer su griglia assoluta & \textbf{2000/2000} & \textbf{9999\us} & \textbf{10\us} & s\`i \\
\bottomrule
\end{tabular}
\caption{Effetto della modalit\`a di attivazione sul task critico, a parit\`a di carico.}
\label{tab:timer}
\end{table}

\textbf{I task di disturbo usano invece l'attesa relativa}, anch'essa per scelta. Essi sono
volutamente in sovraccarico, e con un'attivazione su griglia assoluta un task in ritardo non
attenderebbe pi\`u: il disturbo degenererebbe in un ciclo di calcolo continuo anzich\'e mantenere
il ciclo di lavoro previsto. Sui task di disturbo, del resto, non si misurano scadenze.

\subsection{Le opzioni che governano la strumentazione}

Il comportamento della telemetria \`e determinato da cinque parametri fissati \textbf{al momento
della compilazione}, non modificabili a esecuzione avviata. Ne consegue che ogni configurazione
sperimentale corrisponde a un eseguibile distinto, prodotto passando i relativi valori al
compilatore. La Tabella~\ref{tab:flag} ne riassume il significato.

\begin{table}[h]\centering\small
\begin{tabular}{p{4.6cm}p{2.2cm}p{7.4cm}}
\toprule
\textbf{parametro} & \textbf{valori} & \textbf{significato} \\
\midrule
livello di tracciamento &
0, 1, 2, 3 &
Quanto \`e fine la granularit\`a: 0 disattiva; 1 registra l'esecuzione complessiva e un elemento
per thread; 2 aggiunge un elemento per fase; 3 aggiunge un elemento per ogni iterazione del ciclo,
cio\`e il volume massimo. \\
\addlinespace
tipo di \emph{processor} &
Batch, Simple &
\emph{Come} gli elementi raccolti vengono consegnati. Il primo li accumula in una coda e li invia a
intervalli da un thread dedicato; il secondo li invia \textbf{immediatamente, nel thread stesso che
li ha generati}. \`E il parametro pi\`u rilevante fra quelli studiati. \\
\addlinespace
tipo di \emph{sampler} &
sempre, per proporzione, mai &
\emph{Se} una traccia viene registrata. Nel caso intermedio la decisione \`e probabilistica. \\
\addlinespace
proporzione di campionamento &
$0\ldots1$ &
La probabilit\`a usata dal caso intermedio. \\
\addlinespace
tipo di \emph{exporter} &
collector, uscita standard &
\emph{Dove} gli elementi vengono inviati. \\
\bottomrule
\end{tabular}
\caption{I parametri che governano la strumentazione e ci\`o che ciascuno controlla.}
\label{tab:flag}
\end{table}

La distinzione fra i tre ruoli \`e importante per interpretare i risultati: il \emph{sampler}
decide \emph{se}, il \emph{processor} decide \emph{come} e \emph{quando}, l'\emph{exporter} decide
\emph{dove}. Come si vedr\`a, il costo per il task critico dipende quasi interamente dal secondo.

Sul quinto parametro va fatta una precisazione metodologica. Nella configurazione predefinita gli
elementi vengono inviati a un servizio di raccolta esterno; in assenza di tale servizio l'invio
fallisce senza bloccare l'esecuzione, ma non lascia alcuna traccia contabile di quanti elementi
siano stati effettivamente prodotti. Poich\'e uno degli esperimenti richiede proprio di contarli,
abbiamo reso selezionabile una modalit\`a alternativa che scrive gli elementi sull'uscita standard,
separata dal log del generatore di carico che viene invece emesso sull'uscita di errore. Questa
aggiunta \`e stata realizzata rispettando i parametri esistenti, in modo che le due modalit\`a
differiscano \emph{soltanto} per la destinazione e restino confrontabili sotto ogni altro aspetto.
